\section*{Chapter \ref{ch:b3:advanced-nmr} --- NMR in Depth: Pulses, Relaxation and 2D}

\begin{solution}{exo:b3:advanced-nmr:1}
$\nu_0 = (\gamma/2\pi)B_0$: \ce{^1H} \qty{600.4}{MHz}, \ce{^{13}C} \qty{151.0}{MHz}, \ce{^{19}F}
\qty{565.1}{MHz}.
\end{solution}

\begin{solution}{exo:b3:advanced-nmr:2}
$h\nu_0/2kT = 6.626\times10^{-34} \times 400.2\times10^6/(2 \times 1.381\times10^{-23} \times 300)
= \num{3.2e-5}$. At \qty{4}{K} it is 75 times larger, \num{2.4e-3} (the approximation
still holds).
\end{solution}

\begin{solution}{exo:b3:advanced-nmr:3}
$\theta = \gamma B_1t_p = \pi/2$ in \qty{10}{\micro s}: $\gamma B_1/2\pi = 1/(4 \times
\qty{10}{\micro s}) = \qty{25}{kHz}$. A $180^\circ$ pulse lasts \qty{20}{\micro s}; \qty{5}{\micro s}
gives $45^\circ$.
\end{solution}

\begin{solution}{exo:b3:advanced-nmr:4}
$T_2 = 1/(\pi \times 3.2) = \qty{0.10}{s}$.
\end{solution}

\begin{solution}{exo:b3:advanced-nmr:5}
$(10.71/42.58)^3 = 0.0159$, times 0.011: \num{1.75e-4}, about 1/5700 of the proton
signal. Since S/N $\propto\sqrt N$, equal S/N needs $5700^2 \approx \num{3e7}$ times more
scans --- impossible, which is why ${}^{13}$C spectra use more concentrated samples,
decoupling and polarisation transfer.
\end{solution}

\begin{solution}{exo:b3:advanced-nmr:6}
$T_1 = 1.39/\ln2 = \qty{2.0}{s}$; recycle delay $\ge 5T_1 = \qty{10}{s}$.
\end{solution}

\begin{solution}{exo:b3:advanced-nmr:7}
\omterm{def:b3:advanced-nmr:experiments}{COSY}: one \omterm{def:b3:advanced-nmr:two-d}{cross peak}, between the \ce{CH2} quartet and the \ce{CH3} triplet of the
ethyl group; the \ce{CH3CO} singlet couples to nothing. \omterm{def:b3:advanced-nmr:experiments}{HMBC} from the singlet methyl:
the carbonyl carbon (two bonds) and the \ce{CH2} carbon (three bonds).
\end{solution}

\begin{solution}{exo:b3:advanced-nmr:8}
$k_c = \pi \times 50/\sqrt2 = \qty{111}{s^{-1}}$. $\Delta G^\ddagger = 8.314 \times 330 \times
\ln(6.88\times10^{12}/111) = \qty{68}{kJ/mol}$.
\end{solution}

\begin{solution}{exo:b3:advanced-nmr:9}
NOE $\propto r^{-6}$: $r_b = 0.30 \times 4^{-1/6} = \qty{0.24}{nm}$.
\end{solution}

\begin{solution}{exo:b3:advanced-nmr:10}
In the \omterm{def:b3:advanced-nmr:magnetisation}{rotating frame} on resonance only $\mathbf B_1 = B_1\mathbf e_{x'}$ remains:
$\dot M_{x'} = 0$, $\dot M_{y'} = \gamma B_1M_z$, $\dot M_z = -\gamma B_1M_{y'}$ (components of
$\gamma\mathbf M\times\mathbf B_1$). So $M_{x'}$ is constant and $(M_{y'}, M_z)$ turns at the angular
rate $\gamma B_1$: after $t_p$, by the angle $\gamma B_1t_p$ about $x'$.
\end{solution}

\begin{solution}{exo:b3:advanced-nmr:11}
A spin in a slightly stronger field gains phase at a constant extra rate during
$\tau$; the $180^\circ$ pulse reverses its phase, and it loses exactly the same amount
during the second $\tau$: all such spins realign at $2\tau$. Random spin--spin
interactions fluctuate in time and do not repeat in the second interval: their
dephasing is not undone. The echo amplitude decays with the true $T_2$.
\end{solution}

\begin{solution}{exo:b3:advanced-nmr:12}
Ethyl ethanoate, \ce{CH3COOCH2CH3}: \ce{OCH2} quartet (4.12), \ce{CH3CO} singlet (2.05),
\ce{CH3} triplet; the quartet protons see the carbonyl carbon three bonds away,
through the ester oxygen. Methyl propanoate, \ce{CH3CH2COOCH3}, would show its
quartet near 2.3 ppm (on a carbon near 27 ppm in \omterm{def:b3:advanced-nmr:experiments}{HSQC}, not 60) and a singlet
\ce{OCH3} near 3.7 ppm whose \omterm{def:b3:advanced-nmr:experiments}{HMBC} peak to the carbonyl crosses the oxygen.
\end{solution}

\begin{solution}{pb:b3:advanced-nmr:1}
\textbf{1.} $(2 \times 6 + 2 - 12)/2 = 1$: one C=O.
\textbf{2.} 173.6 ppm: the ester carbonyl; 60.1 ppm: the \ce{OCH2}.
\textbf{3.} Five; DEPT-135: \ce{CH2} negative (60.1, 36.3, 18.6), \ce{CH3} positive (14.3,
13.7), the carbonyl absent.
\textbf{4.} \qty{400.2}{MHz} and \qty{100.7}{MHz}.
\textbf{5.} $3J = \qty{21.6}{Hz} = \qty{0.054}{ppm}$.
\textbf{6.} The multiplicities show which groups are neighbours inside each
fragment, but not how the fragments join across the oxygen and the carbonyl.
\textbf{7.} 4.13--1.25; 2.28--1.66; 1.66--0.95 (and the symmetric ones).
\textbf{8.} \ce{OCH2CH3} (4.13, 1.25) and \ce{CH2CH2CH3} (2.28, 1.66, 0.95).
\textbf{9.} The \ce{OCH2} and \ce{CH2C=O} protons are separated by the oxygen and the
carbonyl carbon: five bonds, no resolved coupling.
\textbf{10.} 1.66 ppm, between \ce{CH2} and \ce{CH3}: coupled to 2 + 3 = 5 protons with
similar $J$, a sextet (a multiplet).
\textbf{11.} Coupled methyls, hence two \ce{CH3} on adjacent carbons (a \ce{CH3-CH3} unit): impossible here.
\textbf{12.} \ce{-O-CH2-CH3} and \ce{-CH2-CH2-CH3}, plus the \ce{C=O}.
\textbf{13.} 4.13/60.1; 2.28/36.3; 1.66/18.6; 1.25/14.3; 0.95/13.7.
\textbf{14.} 14.3 (ethyl \ce{CH3}) and 13.7 ppm (butanoyl \ce{CH3}), only \qty{0.6}{ppm} apart.
\textbf{15.} 4.13 ppm to 173.6 ppm: H--C--O--C, three bonds.
\textbf{16.} 2.28 ppm (two bonds) and 1.66 ppm (three bonds) to 173.6 ppm.
\textbf{17.} \ce{CH3CH2CH2C(=O)OCH2CH3}, ethyl butanoate. Propyl propanoate would put an
\ce{OCH2} triplet (not a quartet) near 4.0 ppm and a quartet near 2.3 ppm.
\textbf{18.} Four-bond couplings are usually below 1 Hz; the \omterm{def:b3:advanced-nmr:experiments}{HMBC} delay, tuned to
5--10 Hz, does not select them.
\textbf{19.} Carbons relax at different rates and receive different nuclear
Overhauser enhancements from decoupling.
\textbf{20.} $1 - \eu^{-5/20} = 0.22$: the carbonyl is under-represented by a factor 4.
\textbf{21.} $1 - \eu^{-5} = 0.993$.
\textbf{22.} The NOE from proton decoupling, different for each carbon; it is removed
by inverse-gated decoupling (decoupler on only during acquisition).
\textbf{23.} $256 \times \qty{100}{s} = \qty{25600}{s}$, about 7 hours.
\textbf{24.} \textbf{$5T_1$ of the carbonyl carbon: \qty{100}{s} between scans.}
\end{solution}
